% Appendix E: further reading. The basics first, then the protocol reference, the standards, the
% tools and four directions further into the subject.

\chapter{Further reading}
\label{app:reading}

This book ends with a CAN node a microcontroller can talk to, and it ends there deliberately: the
design is simplified at several deliberate points (\secref{c19:sec:gaps}), and the driver on the
other side of the SPI link is written by somebody else. This appendix says where the basics the book
assumes can be revised, where the protocol itself is described, which standards the design relates
to, and where to go next.

% ------------------------------------------------------------------------------------------
\section{The basics}
\label{app:reading:basics}

The book assumes the basics of digital electronics and revises them only where the language calls
for it. \debook{} (\url{https://github.com/qrtech-academy/digital-electronics}) is the book to turn
to when they need brushing up. It is freely available, and it ends roughly where this book begins:
with a state machine built from gates and flip-flops, simulated in the same CircuitVerse.

Six chapters plus appendices, and each chapter corresponds to something this book takes for granted:

\begin{booktable}{@{}c>{\raggedright\arraybackslash}p{0.44\linewidth}L@{}}
  \thead{Ch.} & \thead{Subject} & \thead{Here in this book} \\
  \midrule
  1 & Logic gates, and the transistors and the delay inside them & \secref{c1:sec:gates},
    \secref{c5:sec:fmax} \\
  2 & Boolean algebra, De Morgan, number systems and codes & \secref{c1:sec:boolean} \\
  3 & Karnaugh maps and minimisation & \secref{c2:sec:kmap}, \secref{c8:sec:byhand} \\
  4 & Combinational building blocks: multiplexer, decoder, comparator, adder & \secref{c2:sec:mux} \\
  5 & Flip-flops, registers, metastability and synchronisation & chapters~\ref{c3:ch} and~\ref{c4:ch} \\
  6 & Counters, timers and state machines & chapters~\ref{c6:ch}, \ref{c7:ch} and~\ref{c8:ch} \\
  A & Answers to the exercises & \textemdash{} \\
  B & CircuitVerse & \secref{c1:sec:cv} \\
\end{booktable}

The appendices after B are two written papers with model answers. They are the quickest way to test
whether the prerequisites are in place before \chapref{c1:ch}, and sitting one of them against the
clock shows more clearly than a read-through which of the chapters above are worth revising.

% ------------------------------------------------------------------------------------------
\section{The protocol reference}
\label{app:reading:protocol}

\canbook{} is the book this one leans on every time CAN as a \emph{protocol} comes up. It is short
and written for both halves of a CAN system: the one building the controller in hardware and the one
writing the driver in software. It therefore mentions neither course nor class, and it is
deliberately independent of which of the two courses is reading it.

It exists in two editions with the same chapter and section numbering, so every reference in this
book holds for both: in English at
\url{https://github.com/Yrgo-26/can-book/blob/main/en/can-en.pdf}, and in Swedish, as
\textswedish{\emph{CAN: bussen, framen och kontrollern}}, at
\url{https://github.com/Yrgo-26/can-book/blob/main/sv/can-sv.pdf}.

Seven chapters plus an appendix, and the chapter division is this book's part~\ref{part:can} seen from
the protocol side:

\begin{booktable}{@{}c>{\raggedright\arraybackslash}p{0.44\linewidth}L@{}}
  \thead{Ch.} & \thead{Subject} & \thead{Here in this book} \\
  \midrule
  1 & The bus: topology, dominant and recessive, wired-AND, termination & \secref{c10:sec:can} \\
  2 & The frame: the fields, the identifier, RTR, DLC, the acknowledgement &
    \secref{c10:sec:frameformat} \\
  3 & Bit stuffing and CRC-15 & \secref{c10:sec:stuffing}, \chapref{c13:ch}, \ref{c14:ch} and
    \ref{c15:ch} \\
  4 & Arbitration, bit by bit & \secref{c10:sec:arbitration}, \secref{c17:sec:arbitration} \\
  5 & Bit timing and synchronisation & \secref{c10:sec:bittiming}, \secref{c12:sec:bittimer} \\
  6 & Error handling: error frames, error counters, bus-off & \secref{c19:sec:gaps}, as a gap \\
  7 & The simplified controller, and the three limitations that reach all the way up &
    \secref{c19:sec:gaps}, \secref{app:project:limits}, appendix~\ref{app:protocol} \\
  A & Answers to every exercise & \textemdash{} \\
\end{booktable}

Three things make it worth having open beside this book rather than read once:

\begin{itemize}
  \item \textbf{Chapters 1--6 are the requirement specification, bit by bit,} for the controller
    part~\ref{part:can} builds. Where this book says ``the rule is five identical bits in a row'',
    that book says why the rule exists and what breaks without it.
  \item \textbf{Chapter 7 describes exactly this design}, from outside: the blocks, what it
    implements of the standard, what it deliberately leaves out, why a register layer is needed, and
    the three limitations that reach all the way up into the driver's code. Anyone who wants the whole
    picture in five pages before \chapref{c11:ch}'s architecture reads that chapter first.
  \item \textbf{Four exercises per chapter, all pen and paper, all answered in appendix A.} They test
    the protocol understanding without a line of VHDL being involved, which makes them the best
    preparation there is for \chapref{c20:ch}'s reading tasks.
\end{itemize}

The books share an author and are typeset the same way, so a reference looks the same in both. They do
number their chapters independently of each other, however: ``chapter~3'' with no further
qualification always means a chapter in \emph{this} book, and a reference to the other is always
written out in full.

% ------------------------------------------------------------------------------------------
\section{Standards and specifications}
\label{app:reading:specs}

The book relates to real CAN throughout without being it, and \canbook{} is the intermediate step:
it describes the standard at the level of detail a designer needs. If you then want to see what the
original actually says:

\begin{itemize}
  \item \textbf{Bosch CAN Specification 2.0} (parts A and B) is the base document: the frame formats,
    the bit stuffing, the CRC polynomial and the arbitration. Part A is the standard frame with an
    11-bit identifier, that is, this book's frame; part B adds the extended 29-bit form.
  \item \textbf{ISO 11898-1} is the standardisation of the same thing, with the data link layer and
    the physical signalling interface. ISO 11898-2 is the high-speed transceiver, that is, the device
    that would sit between \code{tx_bus}/\code{bus_en} and a real differential bus.
  \item \textbf{IEEE 1076-1993} is VHDL-93, the language standard \code{--std=93} selects. It is worth
    reading as a reference work rather than cover to cover; the packages \code{ieee.std_logic_1164}
    and \code{ieee.numeric_std} are the two parts you have used throughout the book.
\end{itemize}

The four points where the book's controller deliberately departs from the standard are collected in
\secref{c19:sec:gaps}, and \secref{app:project:limits} says which three of them reach all the way out
into the register map.

% ------------------------------------------------------------------------------------------
\section{The tools}
\label{app:reading:tools}

\begin{itemize}
  \item \textbf{GHDL's documentation} covers considerably more than the three commands
    appendix~\ref{app:sim} needs: waveform output to VCD or GHW, code coverage, and the other
    language standard flags. \textbf{GTKWave} is the usual waveform viewer to pair it with, and the
    first time you debug a timing bug by \emph{looking} at the signals instead of reading
    \code{report} lines it is worth its installation time.
  \item \textbf{Quartus Prime Lite's own documentation}, in particular the timing analysis:
    appendix~\ref{app:quartus} settles for one \code{create_clock}, while a real design also needs
    input and output constraints and exceptions for crossing clock domains.
  \item \textbf{The DE0-CV's user manual} from the board manufacturer has the complete pin table, the
    peripherals and the schematic. It is the answer key when \file{DE0_CV.qsf} does not cover
    something you want to use.
  \item \textbf{CircuitVerse} is free and runs in the browser. It is still the fastest way to try a
    small combinational idea, and \chapref{c9:ch} and \ref{c20:ch} allow it for exactly that reason.
\end{itemize}

% ------------------------------------------------------------------------------------------
\section{Further into digital design}
\label{app:reading:design}

This is a book about one system, not about the whole subject. Four directions that lie close by and
for which material is not lacking:

\begin{itemize}
  \item \textbf{Timing analysis and clock domains in earnest.} \Chapref{c4:ch} gives the
    two-flip-flop synchroniser and \secref{c5:sec:fmax} a first Fmax. The next step is setup and hold
    analysis, FIFOs between clock domains, and Gray-coded pointers \textemdash{} that is, what you do
    when a whole bus, and not a single bit, has to pass between two clocks.
  \item \textbf{Verification beyond a hand-written testbench:} self-checking models, random stimuli
    and coverage metrics. VHDL's OSVVM and UVVM are the two common frameworks, and the step from
    there to SystemVerilog and UVM is the one industry takes.
  \item \textbf{Other interfaces, with the same method.} The book builds CAN; UART, SPI,
    I\textsuperscript{2}C and Modbus all have the same shape \textemdash{} a bit timer, a shift
    register, a sequencer \textemdash{} and building one more of them with the book's patterns is the
    fastest way to notice how much of \chapref{c11:ch}'s split was common property.
  \item \textbf{Processor design.} A small RISC core is the natural next step for anyone who found the
    state machine in \chapref{c16:ch} the most interesting part: the same building blocks, an
    instruction decoder instead of a frame decoder, and a register file interface instead of a bus.
\end{itemize}

What all four have in common is that they rest on the same five patterns \chapref{c20:ch} lists. That
is the real takeaway from this book: not CAN, but that you can build the next protocol too.
